\subsection{A PRELIMINARY CONSTRUCTION}

In this number and the next, we denote by R a reduced root system in a vector space V and by B a basis of R. We denote by \((n(\alpha,\beta))_{\alpha,\beta\in\mathbb{B}}\) the Cartan matrix relative to B. Recall that \(n(\alpha,\beta)=\langle\alpha,\beta^{\vee}\rangle\) . We are going to show that R is the root system of a split semi-simple Lie algebra which is unique up to isomorphism. In the main we shall be considering the Lie algebra defined by the relations in Prop. 1.

The construction in this number applies to any square matrix \((n(\alpha,\beta))_{\alpha,\beta\in\mathbb{B}}\) over k with non-zero determinant and such that \(n(\alpha,\alpha)=2\) for all \(\alpha\in B\) (cf.~Chap. VI, §1, no. 10, formula (14)).

Let E be the free associative algebra of the set B over k. Recall that E is N-graded (Algebra, Chap. III, §3, no. 1, Example 3). We are going to associate to each \(\alpha\in B\) endomorphisms \(X_{-\alpha}^{0}, H_{\alpha}^{0}, X_{\alpha}^{0}\) of the vector space E, of degrees 1, 0, -1 respectively. For any word \((\alpha_{1},\ldots,\alpha_{n})\) in elements of B, put

\[
X _ {- \alpha} ^ {0} (\alpha_ {1}, \ldots , \alpha_ {n}) = (\alpha , \alpha_ {1}, \ldots , \alpha_ {n})\tag{7}
\]

\[
H _ {\alpha} ^ {0} \left(\alpha_ {1}, \dots , \alpha_ {n}\right) = \left(- \sum_ {i = 1} ^ {n} n \left(\alpha_ {i}, \alpha\right)\right) \left(\alpha_ {1}, \dots , \alpha_ {n}\right).\tag{8}
\]

On the other hand, \(X_{\alpha}^{0}(\alpha_{1},\ldots ,\alpha_{n})\) is defined by induction on \(n\) using the formula

\[
X _ {\alpha} ^ {0} (\alpha_ {1}, \ldots , \alpha_ {n}) = (X _ {- \alpha_ {1}} ^ {0} X _ {\alpha} ^ {0} - \delta_ {\alpha , \alpha_ {1}} H _ {\alpha} ^ {0}) (\alpha_ {2}, \ldots , \alpha_ {n})\tag{9}
\]

where \(\delta_{\alpha, \alpha_1}\) is the Kronecker symbol; it is understood that \(X_{\alpha}^{0}(\alpha_{1}, \ldots, \alpha_{n})\) is zero if \((\alpha_{1}, \ldots, \alpha_{n})\) is the empty word.

Lemma 1. For all \(\alpha, \beta \in \mathbf{B}\) , we have

\[
[ X _ {\alpha} ^ {0}, X _ {- \alpha} ^ {0} ] = - H _ {\alpha} ^ {0}\tag{10}
\]

\[
[ H _ {\alpha} ^ {0}, H _ {\beta} ^ {0} ] = 0\tag{11}
\]

\[
[ H _ {\alpha} ^ {0}, X _ {\beta} ^ {0} ] = n (\beta , \alpha) X _ {\beta} ^ {0}\tag{12}
\]

\[
[ H _ {\alpha} ^ {0}, X _ {- \beta} ^ {0} ] = - n (\beta , \alpha) X _ {- \beta} ^ {0}\tag{13}
\]

\[
[ X _ {\alpha} ^ {0}, X _ {- \beta} ^ {0} ] = 0 \text {   if   } \alpha \neq \beta .\tag{14}
\]

Indeed, relation (9) can be written

\[
(X _ {\alpha} ^ {0} X _ {- \alpha_ {1}} ^ {0}) (\alpha_ {2}, \ldots , \alpha_ {n}) = (X _ {- \alpha_ {1}} ^ {0} X _ {\alpha} ^ {0}) (\alpha_ {2}, \ldots , \alpha_ {n}) - \delta_ {\alpha , \alpha_ {1}} H _ {\alpha} ^ {0} (\alpha_ {2}, \ldots , \alpha_ {n})
\]

which proves (10) and (14). Relation (11) is clear. Next

\[
\begin{array}{r l} [ H _ {\alpha} ^ {0}, X _ {- \beta} ^ {0} ] (\alpha_ {1}, \ldots , \alpha_ {n}) & = H _ {\alpha} ^ {0} (\beta , \alpha_ {1}, \ldots , \alpha_ {n}) + \left(\sum_ {i = 1} ^ {n} n (\alpha_ {i}, \alpha)\right) (\beta , \alpha_ {1}, \ldots , \alpha_ {n}) \\ & = - n (\beta , \alpha) (\beta , \alpha_ {1}, \ldots , \alpha_ {n}) \\ & = - n (\beta , \alpha) X _ {- \beta} ^ {0} (\alpha_ {1}, \ldots , \alpha_ {n}) \end{array}
\]

hence (13). Finally,

\[
\begin{array}{l} 0 = [ H _ {\alpha} ^ {0}, [ X _ {\beta} ^ {0}, X _ {- \gamma} ^ {0} ] ] \quad \text { by (10), (11), (14)} \\ \qquad = [ [ H _ {\alpha} ^ {0}, X _ {\beta} ^ {0} ], X _ {- \gamma} ^ {0} ] + [ X _ {\beta} ^ {0}, [ H _ {\alpha} ^ {0}, X _ {- \gamma} ^ {0} ] ] \\ \qquad = [ [ H _ {\alpha} ^ {0}, X _ {\beta} ^ {0} ] - n (\gamma , \alpha) X _ {\beta} ^ {0}, X _ {- \gamma} ^ {0} ] \quad \text { by (13)} \\ \qquad = [ [ H _ {\alpha} ^ {0}, X _ {\beta} ^ {0} ] - n (\beta , \alpha) X _ {\beta} ^ {0}, X _ {- \gamma} ^ {0} ] \quad \text { by (14); } \end{array}\tag{15}
\]

now, considering the empty word immediately gives

\[
([ H _ {\alpha} ^ {0}, X _ {\beta} ^ {0} ] - n (\beta , \alpha) X _ {\beta} ^ {0}) (\emptyset) = 0
\]

so (15) implies that

\[
\left(\left[ H _ {\alpha} ^ {0}, X _ {\beta} ^ {0} \right] - n (\beta , \alpha) X _ {\beta} ^ {0}\right) X _ {- \gamma_ {1}} ^ {0} X _ {- \gamma_ {2}} ^ {0} \dots X _ {- \gamma_ {n}} ^ {0} (\varnothing) = 0
\]

for all \(\gamma_1, \ldots, \gamma_n \in \mathrm{B}\) ; this proves (12).

Lemma 2. The endomorphisms \(X_{\alpha}^{0}\) , \(H_{\beta}^{0}\) , \(X_{-\gamma}^{0}\) , where \(\alpha, \beta, \gamma \in \mathrm{B}\) , are linearly independent.

Since \(X_{-\alpha}^{0}(\varnothing) = \alpha\) , it is clear that the \(X_{-\alpha}^{0}\) are linearly independent. Assume that \(\sum_{\alpha} a_{\alpha} H_{\alpha}^{0} = 0\) ; then, for all \(\beta \in B\) ,

\[
0 = \left[ \sum_ {\alpha} a _ {\alpha} H _ {\alpha} ^ {0}, X _ {- \beta} ^ {0} \right] = - \sum_ {\alpha} a _ {\alpha} n (\beta , \alpha) X _ {- \beta} ^ {0};
\]

since \(\operatorname{det}(n(\beta, \alpha)) \neq 0\) , it follows that \(a_{\alpha} = 0\) for all \(\alpha\) . Assume that \(\sum_{\alpha} a_{\alpha} X_{\alpha}^{0} = 0\) . In view of formulas (7), (8), (9),

\[
X _ {\alpha} ^ {0} (\beta) = 0,
\]

\[
X _ {\alpha} ^ {0} (\beta , \beta) = 2 \delta_ {\alpha \beta} \beta
\]

for all \(\beta \in \mathrm{B}\) . It follows that \(a_{\beta} = 0\) for all \(\beta\) . Since \(X_{\alpha}^{0}\) , \(H_{\alpha}^{0}\) , \(X_{-\alpha}^{0}\) are of degree \(-1, 0, 1\) , respectively, the lemma follows from what has gone before.

Let I be the set \(B \times \{-1,0,1\}\) . Put \(x_{\alpha} = (\alpha, -1)\) , \(h_{\alpha} = (\alpha, 0)\) , and \(x_{-\alpha} = (\alpha, 1)\) . Let a be the Lie algebra defined by the generating family I and the following set R of relators:

\[
\begin{array}{l} \left[ h _ {\alpha}, h _ {\beta} \right] \\ \left[ h _ {\alpha}, x _ {\beta} \right] - n (\beta , \alpha) x _ {\beta} \\ \left[ h _ {\alpha}, x _ {- \beta} \right] + n (\beta , \alpha) x _ {- \beta} \\ \left[ x _ {\alpha}, x _ {- \alpha} \right] + h _ {\alpha} \\ \left[ x _ {\alpha}, x _ {- \beta} \right] \quad \text {if} \alpha \neq \beta \end{array}
\]

(cf.~Chap. II, §2, no. 3). By Lemma 1, there exists a unique linear representation \(\rho\) of \(\mathfrak{a}\) on \(\mathrm{E}\) such that

\[
\rho (x _ {\alpha}) = X _ {\alpha} ^ {0}, \rho (h _ {\alpha}) = H _ {\alpha} ^ {0}, \rho (x _ {- \alpha}) = X _ {- \alpha} ^ {0}.
\]

In view of Lemma 2, this proves the following result:

Lemma 3. The canonical images in \(\mathfrak{a}\) of the elements \(x_{\alpha}, h_{\beta}, x_{-\gamma}\) , where \(\alpha, \beta, \gamma \in \mathrm{B}\) , are linearly independent.

In the following, we identify \(x_{\alpha}\) , \(h_{\alpha}\) , \(x_{-\alpha}\) with their canonical images in a.

Lemma 4. There exists a unique involutive automorphism \(\theta\) of \(\mathfrak{a}\) such that

\[
\theta (x _ {\alpha}) = x _ {- \alpha}, \theta (x _ {- \alpha}) = x _ {\alpha}, \theta (h _ {\alpha}) = - h _ {\alpha}
\]

for all \(\alpha\in B\) .

Indeed, there exists an involutive automorphism of the free Lie algebra \(L(I)\) satisfying these conditions. It leaves \(R \cup (-R)\) stable, and hence defines by passage to the quotient an involutive automorphism of a satisfying the conditions of the lemma. The uniqueness follows from the fact that a is generated by the elements \(x_{\alpha}\) , \(h_{\alpha}\) , \(x_{-\alpha}\) ( \(\alpha \in B\) ).

This automorphism is called the canonical involutive automorphism of a.

Let Q be the set of radical weights of R; this is a free Z-module with basis B (Chap. VI, §1, no. 9). There exists a graduation of type Q on the free Lie algebra L(I) such that \(x_{\alpha}\) , \(h_{\alpha}\) , \(x_{-\alpha}\) are of degrees \(\alpha\) , 0, - \(\alpha\) , respectively (Chap. II, §2, no. 6). Now the elements of R are homogeneous. Hence there exists a unique graduation of type Q on a compatible with the Lie algebra structure of a and such that \(x_{\alpha}\) , \(h_{\alpha}\) , \(x_{-\alpha}\) are of degrees \(\alpha\) , 0, - \(\alpha\) , respectively. For any \(\mu \in Q\) , denote by \(a^{\mu}\) the set of elements of a homogeneous of degree \(\mu\) .

Lemma 5. Let \(z \in \mathfrak{a}\) . Then \(z \in \mathfrak{a}^{\mu}\) if and only if \([h_{\alpha}, z] = \langle \mu, \alpha^{\vee} \rangle z\) for all \(\alpha \in \mathrm{B}\) .

For \(\mu \in \mathrm{Q}\) , let \(\mathfrak{a}^{(\mu)}\) be the set of \(x \in \mathfrak{a}\) such that \([h_{\alpha}, x] = \langle \mu, \alpha^{\vee} \rangle x\) for all \(\alpha \in \mathrm{B}\) . The sum of the \(\mathfrak{a}^{(\mu)}\) is direct. To prove the lemma, it therefore suffices to show that \(\mathfrak{a}^{\mu} \subset \mathfrak{a}^{(\mu)}\) . Let \(\alpha \in \mathrm{B}\) . The endomorphism \(u\) of the vector space \(\mathfrak{a}\) such that \(u|\mathfrak{a}^{\mu} = \langle \mu, \alpha^{\vee} \rangle\) .1 is a derivation of \(\mathfrak{a}\) such that \(ux = (\operatorname{ad} h_{\alpha})\) . \(x\) for \(x = x_{\beta}\) , \(x = h_{\beta}\) , \(x = x_{-\beta}\) ; hence \(u = \operatorname{ad} h_{\alpha}\) , which proves our assertion.

Remark. It follows from Lemma 5 that every ideal of \(\mathfrak{a}\) is homogeneous, since it is stable under the ad \(h_{\alpha}\) .

Denote by \(Q_{+}\) (resp. \(Q_{-}\) ) the set of linear combinations of elements of B with positive (resp. negative) integer coefficients, not all zero. Put \(a_{+} = \sum_{\mu \in Q_{+}} a^{\mu}\) and \(a_{-} = \sum_{\mu \in Q_{-}} a^{\mu}\) . Since \(Q_{+} + Q_{+} \subset Q_{+}\) and \(Q_{-} + Q_{-} \subset Q_{-}\) , \(a_{+}\) and \(a_{-}\) are Lie subalgebras of a.

PROPOSITION 2. (i) The Lie algebra \(\mathfrak{a}_{+}\) is generated by the family \((x_{\alpha})_{\alpha \in \mathbb{B}}\) .

\begin{enumerate}
\def\labelenumi{(\roman{enumi})}
\setcounter{enumi}{1}
\item
  The Lie algebra \(\mathfrak{a}_{-}\) is generated by the family \((x_{-\alpha})_{\alpha \in \mathrm{B}}\) .
\item
  The family \((h_{\alpha})_{\alpha \in \mathrm{B}}\) is a basis of the vector space \(\mathfrak{a}^0\) .
\item
  The vector space \(\mathfrak{a}\) is the direct sum of \(\mathfrak{a}_{+},\mathfrak{a}^{0},\mathfrak{a}_{-}\) .
\end{enumerate}

Let \(\mathfrak{r}\) (resp. \(\mathfrak{n}\) ) be the Lie subalgebra of \(\mathfrak{a}\) generated by \((x_{\alpha})_{\alpha \in \mathrm{B}}\) (resp. \((x_{-\alpha})_{\alpha \in \mathrm{B}}\) ), and \(\mathfrak{h}\) the vector subspace of \(\mathfrak{a}\) generated by \((h_{\alpha})_{\alpha \in \mathrm{B}}\) . Since the \(x_{\alpha}\) are homogeneous elements of \(\mathfrak{a}_{+}\) , \(\mathfrak{r}\) is a graded subalgebra of \(\mathfrak{a}_{+}\) ; hence, \([\mathfrak{h},\mathfrak{r}]\subset \mathfrak{r}\) , so \(\mathfrak{h} + \mathfrak{r}\) is a subalgebra of \(\mathfrak{a}\) ; since

\[
[ x _ {- \alpha}, x _ {\beta} ] = \delta_ {\alpha \beta} h _ {\alpha},
\]

\([x_{-\alpha}, \mathfrak{r}] \subset \mathfrak{h} + \mathfrak{r}\) for all \(\alpha \in B\) . Similarly, \(\mathfrak{n}\) is a graded subalgebra of \(\mathfrak{a}_{-}\) , one has \([\mathfrak{h}, \mathfrak{n}] \subset \mathfrak{n}\) , \(\mathfrak{h} + \mathfrak{n}\) is a subalgebra of \(\mathfrak{n}\) , and \([x_{\alpha}, \mathfrak{n}] \subset \mathfrak{h} + \mathfrak{n}\) for all \(\alpha \in B\) . Put \(a' = r + h + n\) . The preceding shows that \(a'\) is stable under ad \(x_{\alpha}\) , ad \(h_{\alpha}\) and ad \(x_{-\alpha}\) for all \(\alpha \in B\) , and hence is an ideal of a. Since \(a'\) contains \(x_{\alpha}\) , \(h_{\alpha}\) , \(x_{-\alpha}\) for all \(\alpha \in B\) , \(a' = a\) . It follows from this that the inclusions \(r \subset a_{+}\) , \(h \subset a^{0}\) , \(n \subset a_{-}\) are equalities, which proves the proposition.

PROPOSITION 3. The Lie algebra \(\mathfrak{a}_{+}\) (resp. \(\mathfrak{a}_{-}\) ) is a free Lie algebra with basic family \((x_{\alpha})_{\alpha \in \mathrm{B}}\) (resp. \((x_{-\alpha})_{\alpha \in \mathrm{B}}\) ) (cf.~Chap. II, §2, no. 3).

Let L be the Lie subalgebra of E generated by B. By Chap. II, §3, Th. 1, L can be identified with the free Lie algebra generated by B. The left regular representation of E on itself is clearly injective, and defines by restriction to L an injective representation \(\rho'\) of the Lie algebra L on E. Let \(\varphi\) be the unique homomorphism from L to \(\mathfrak{a}_{-}\) which takes \(\alpha\) to \(x_{-\alpha}\) for all \(\alpha \in B\) . Then, for all \(\alpha \in B\) , \(\rho(\varphi(\alpha))\) is the endomorphism of left multiplication by \(\alpha\) on E, so \(\rho \circ \varphi = \rho'\) , which proves that \(\varphi\) is injective. Thus, \((x_{-\alpha})_{\alpha \in B}\) is a basic family for \(\mathfrak{a}_{-}\) . Since \(\theta(x_{-\alpha}) = x_{\alpha}\) for all \(\alpha\) (cf.~Lemma 4), \((x_{\alpha})_{\alpha \in B}\) is a basic family for \(\mathfrak{a}_{+}\) .
% semantic-fragments: legacy-input-seam
\relax{}
